\section{Experimental Protocol Details}
\label{app:experimental-protocol}

This appendix provides complete experimental methodology details for reproducibility.

\subsection{Software and Hardware}

\textbf{METIS Configuration}:
\begin{itemize}
    \item Version: METIS 5.1.0 (March 2013 release)
    \item Wrapper: Python 3.11 with METIS C library via ctypes
    \item Algorithm: Multilevel k-way partitioning (METIS\_PartGraphKway for n-way), recursive bisection (METIS\_PartGraphRecursive for bisection)
    \item Command-line flags: \texttt{-ufactor=5} (0.5\% imbalance tolerance), \texttt{-niter=10} (10 refinement iterations), \texttt{-seed=42} (reproducible random seed)
    \item Edge weight handling: Symmetric weights (if edge $(u,v)$ has weight $w$, then edge $(v,u)$ has weight $w$)
\end{itemize}

\textbf{Hardware Specifications}:
\begin{itemize}
    \item CPU: AMD Ryzen 9 5950X (16 cores, 3.4 GHz base / 4.9 GHz boost)
    \item RAM: 64 GB DDR4-3600
    \item Storage: 2 TB NVMe SSD (Samsung 980 PRO)
    \item OS: Windows 11 Pro (Build 26200), Python 3.13
\end{itemize}

\textbf{Runtime Environment}:
\begin{itemize}
    \item Single-threaded execution (METIS does not use parallelism by default)
    \item No other compute-intensive processes during experiments
    \item Total wall-clock time for all 43 runs: 47 minutes (mean: 66 seconds per run)
\end{itemize}

\subsection{Data Preprocessing}

\textbf{Census Tract Data}:
\begin{itemize}
    \item Source: US Census Bureau 2020 Redistricting Data (P.L. 94-171)
    \item Downloaded: April 2021 (official release)
    \item Tract count by state: Alabama (1,181), Georgia (1,969), Louisiana (1,148), Mississippi (667), South Carolina (1,103)
    \item Total tracts: 6,068
\end{itemize}

\textbf{Minority VAP Calculation}:
\begin{itemize}
    \item Voting Age Population (VAP) = Population aged 18+
    \item Minority VAP = Black VAP + Hispanic VAP + Asian VAP + Native American VAP + Pacific Islander VAP + Two or More Races VAP
    \item Non-minority (White) VAP = Total VAP - Minority VAP
    \item Minority percentage = Minority VAP / Total VAP
\end{itemize}

\textbf{Adjacency Graph Construction}:
\begin{itemize}
    \item Method: Shared boundary detection using GeoPandas 0.12.2
    \item Two tracts are adjacent if their geometries share a boundary segment (not just a point)
    \item Adjacency is symmetric: if tract $i$ is adjacent 	o tract $j$, then tract $j$ is adjacent 	o tract $i$
    \item Average degree: Alabama: 5.8, Georgia: 5.6, Louisiana: 5.9, Mississippi: 5.7, South Carolina: 5.7
    \item Graph connectivity: All adjacency graphs are connected (single component)
\end{itemize}

\textbf{Edge Weight Assignment}:
\begin{itemize}
    \item Minority threshold: $\tau = 0.40$ (40\% minority VAP)
    \item Edge weight factor: $\alpha = 5.0$
    \item Weight formula: $w(u,v) = \alpha$ if $m_u \geq \tau$ and $m_v \geq \tau$, otherwise $w(u,v) = 1$
    \item Weight symmetry: $w(u,v) = w(v,u)$ for all edges $(u,v)$
    \item Minority-minority edge fraction: Alabama: 12.3\%, Georgia: 18.7\%, Louisiana: 15.1\%, Mississippi: 14.8\%, South Carolina: 11.9\%
\end{itemize}

\subsection{Experimental Design}

\textbf{Test States}:
Five states selected based on:
\begin{enumerate}
    \item High minority population (Black+Hispanic VAP $> 35\%$)
    \item Active VRA Section 2 litigation history
    \item Range of district counts ($k \in \{4, 6, 7, 14\}$) 	o test scalability
    \item Geographic diversity (Deep South + Atlantic coast)
\end{enumerate}

\textbf{Partitioning Methods}:
\begin{enumerate}
    \item \textbf{Predetermined Recursive Bisection}: For each $k$, test all possible first split configurations: $(1, k-1), (2, k-2), \ldots, (\lfloor k/2 \rfloor, \lceil k/2 \rceil)$. Example: $k=7$ tests 3 configurations: $(1,6), (2,5), (3,4)$.

    \item \textbf{Adaptive Recursive Bisection}: Same as predetermined but evaluates all $k$ candidate first splits and selects the one maximizing minority concentration. (In our experiments, adaptive always selected $(1, k-1)$ due 	o edge-weighting.)

    \item \textbf{N-way Partitioning}: Single METIS call with $k$ partitions directly (no recursive splitting).
\end{enumerate}

\textbf{Total Runs}:
\begin{itemize}
    \item Alabama ($k=7$): 3 predetermined + 1 adaptive + 1 n-way = 5 runs
    \item Georgia ($k=14$): 7 predetermined + 1 adaptive + 1 n-way = 9 runs
    \item Louisiana ($k=6$): 3 predetermined + 1 adaptive + 1 n-way = 5 runs
    \item Mississippi ($k=4$): 2 predetermined + 1 adaptive + 1 n-way = 4 runs
    \item South Carolina ($k=7$): 3 predetermined + 1 adaptive + 1 n-way = 5 runs
    \item \textbf{Total}: 28 runs
\end{itemize}

\textbf{Additional Experiments} (Section 6.2):
\begin{itemize}
    \item Robustness test: 10 random seeds for each method 	imes each state = 150 runs
    \item Tree structure experiment: All Catalan number trees for $k \in \{4,5,6\}$ = 65 additional runs
    \item \textbf{Grand Total}: 243 runs
\end{itemize}

\subsection{Randomness and Reproducibility}

\textbf{METIS Random Seed}:
\begin{itemize}
    \item Fixed seed: \texttt{-seed=42} for all runs in main experiment (Section 5)
    \item Multiple seeds: \texttt{-seed=\\{1,2,3,...,10\\}} for robustness test (Section 6.2)
    \item Seed affects: tie-breaking during coarsening, initial partition during refinement
    \item Seed does \textit{not} affect: graph structure, edge weights, balance constraints
\end{itemize}

\textbf{Single Run Policy}:
For the main experiment (Table 2-3, Figures 1-6), we report results from a \textit{single run} per method per state using \texttt{-seed=42}. This is justified because:
\begin{enumerate}
    \item Edge-weighting with $\alpha = 5$ induces deterministic behavior (Theorem 1)
    \item Robustness test (Section 6.2) shows variance across seeds is negligible ($< 0.1\%$)
    \item Single-run policy matches prior graph partitioning literature~\cite{karypis1998fast,karypis2000multilevel}
\end{enumerate}

For the robustness test, we average over 10 seeds and report mean pm standard deviation.

\subsection{Metrics Calculation}

\textbf{Primary Metrics}:
\begin{itemize}
    \item \textbf{Max Minority Percentage}: $\max_{i \in [k]} \frac{M_i}{N_i}$ where $M_i$ is minority VAP in district $i$ and $N_i$ is total VAP in district $i$

    \item \textbf{MM District Count}: Number of districts with minority percentage $\geq 50\%$ (majority-minority districts under VRA)

    \item \textbf{Population Balance}: Maximum deviation from target population. Constraint: $\leq 0.5\%$ (METIS \texttt{-ufactor=5})
\end{itemize}

\textbf{Secondary Metrics} (Section 5.3):
\begin{itemize}
    \item \textbf{Polsby-Popper Score}: $\text{PP}(D) = \frac{4\pi \cdot \text{Area}(D)}{\text{Perimeter}(D)^2}$, computed using GeoPandas \texttt{.area} and \texttt{.length} methods on dissolved tract geometries

    \item \textbf{Objective Value}: Weighted edge-cut = $\sum_{(u,v) \in \text{cut edges}} w(u,v)$, where cut edge means $u$ and $v$ are in different districts

    \item \textbf{Runtime}: Wall-clock time from start of METIS call 	o return, measured using Python \texttt{time.perf\_counter()}
\end{itemize}

\subsection{Statistical Analysis}

\textbf{Variance Calculation}:
For each state, compute variance across methods:
\[
\text{Var}(X) = \frac{1}{n-1} \sum_{i=1}^{n} (x_i - \bar{x})^2
\]
where $n$ is number of methods (5-9 depending on state) and $x_i$ is max minority percentage for method $i$.

\textbf{Significance Testing}:
We do not perform statistical significance tests (e.g., t-tests, ANOVA) because:
\begin{enumerate}
    \item Sample size is small ($n \leq 9$ methods)
    \item Observed variance is essentially zero ($< 0.1\%$), making significance testing unnecessary
    \item Determinism is theoretical (Theorem 1), not empirical statistical inference
\end{enumerate}

\textbf{Reporting Precision}:
\begin{itemize}
    \item Minority percentages: 1 decimal place (e.g., 66.2\%)
    \item Compactness scores: 2 decimal places (e.g., 0.42)
    \item Runtimes: integer seconds (e.g., 83 seconds)
    \item Variances: 4 decimal places (e.g., 0.0001)
\end{itemize}

\subsection{Code Availability}

All code, data, and results are available at:
\begin{itemize}
    \item GitHub repository: \texttt{https://github.com/[anonymized-for-review]}
    \item Includes: Python scripts, METIS wrapper, tract data processing, visualization code
    \item License: MIT (open source)
    \item Documentation: README with step-by-step reproduction instructions
\end{itemize}

\textbf{Reproduction Steps}:
\begin{enumerate}
    \item Download census tract shapefiles and redistricting data
    \item Run \texttt{build\_adjacency.py} 	o construct graphs
    \item Run \texttt{compute\_edge\_weights.py} with $\alpha=5$, $\tau=0.40$
    \item Run \texttt{run\_partitioning.py} for each method and state
    \item Run \texttt{analyze\_results.py} 	o compute metrics
    \item Run \texttt{create\_visualizations.py} 	o generate figures
\end{enumerate}

Expected runtime: $< 2$ hours on modern hardware.
